\section{Battery integration}
For calculating the absorbed/desorbed amount of energy in a battery, I had to integrate over the domain. I found an easier way to tackle this problem by using the mass matrix and fem indeces already calculated in the spacial discretization. Below is a quick description of this method, which I could not really find a logical place in my thesis for, apart from the appendix.

The amount of absorbed gas stored in the solid phase is obtained by integrating the solid density over the domain:
\begin{equation}
M_s(t)=\int_{\Omega}\rho_s(x,t)\,dx .
\end{equation}
The finite element approximation of the solid density is
\begin{equation}
\rho_s^n(x)=\sum_{i=1}^{N_s} \rho_s^{n,i}\phi_i(x),
\end{equation}
where $N_s$ is the number of solid density degrees of freedom and
$\phi_i$ are the finite element basis functions. Therefore,
\begin{equation}
M_s^n
=
\int_{\Omega}
\sum_{i=1}^{N_s}\rho_s^{n,i}\phi_i(x)\,dx .
\end{equation}
Rearranging the summation gives
\begin{equation}
M_s^n
=
\sum_{i=1}^{N_s}
\rho_s^{n,i}
\int_{\Omega}\phi_i(x)\,dx .
\end{equation}
The integral of each basis function is obtained from the consistent mass matrix.
The mass matrix is defined as
\begin{equation}
(M_s)_{ij}
=
\int_{\Omega}\phi_i(x)\phi_j(x)\,dx .
\end{equation}
Multiplying the mass matrix by a vector of ones gives
\begin{equation}
b_i
=
\sum_{j=1}^{N_s}(M_s)_{ij}
=
\sum_{j=1}^{N_s}
\int_{\Omega}\phi_i(x)\phi_j(x)\,dx .
\end{equation}
Since the finite element basis satisfies the partition of unity,
\begin{equation}
\sum_j\phi_j(x)=1,
\end{equation}
this becomes
\begin{equation}
b_i
=
\int_{\Omega}\phi_i(x)\,dx .
\end{equation}
The total absorbed mass at timestep $n$ is therefore approximated by
\begin{equation}
M_s^n
\approx
\sum_{i=1}^{N_s}b_i\rho_s^{n,i}
=
\mathbf{b}^T\boldsymbol{\rho}_s^n .
\end{equation}  
The maximum possible absorbed mass corresponds to a completely saturated
solid phase:
\begin{equation}
M_{\mathrm{sat}}
=
\int_{\Omega}\rho_{\mathrm{sat}}\,dx
=
\rho_{\mathrm{sat}}|\Omega|,
\end{equation}
where $|\Omega|$ is the domain length (equal to $1$ in the current model).
The filling fraction of the solid storage capacity is then defined as
\begin{equation}
f^n
=
\frac{M_s^n}{M_{\mathrm{sat}}}
=
\frac{\mathbf{b}^T\boldsymbol{\rho}_s^n}
{\rho_{\mathrm{sat}}|\Omega|}.
\end{equation}
A value of
\begin{equation}
f^n=0
\end{equation}
corresponds to an empty solid phase, while
\begin{equation}
f^n=1
\end{equation}
corresponds to a completely saturated solid phase.

\newpage